\section{相关技术与理论基础}

这一章我们梳理了本课程设计用到的五个方向的基础知识：视觉语言模型、CLIP 语义空间、EEG 视觉解码、视觉退化建模与多模态融合，还有 LoRA 参数高效微调。每小节只讲和我们设计直接挂钩的核心思路和方法。

\subsection{视觉语言模型概述（VLM）}

视觉语言模型说白了就是能同时看图和读字、在两种模态之间做信息融合和推理的深度学习模型。按输出形式可以分为两大类：对齐型和生成型。

对齐型 VLM 最出名的是 CLIP。CLIP 用大概 4 亿对图文数据做对比预训练，让图像编码器（ViT 或者 ResNet）和文本编码器（Transformer）学到一个共享的对齐嵌入空间。训练收敛之后，语义上匹配的图和文在向量空间里距离就近，不匹配的自然就远，所以直接算余弦相似度就能做零样本分类。对齐型 VLM 产出的是一定维度的嵌入向量，适合拿来做检索、分类和特征提取这类活儿。

生成型 VLM 在对齐型的基础上又加了一个语言解码模块。典型的有 BLIP-2、LLaVA 还有 Qwen-VL 系列。BLIP-2 搞了个 Q-Former，把视觉特征压成固定数量的查询 token，然后塞进冻结的大语言模型里；LLaVA 更简单，就一个线性投影层把视觉 token 映射到语言模型的嵌入空间完事；Qwen-VL 是阿里通义千问团队出的原生态多模态模型，视觉和语言模块在预训练阶段就已经做了端到端的图文对齐。

这两类模型的根本差别在于输出形式、训练目标和适用场景。对齐型 VLM 只有编码器，输出固定维度的向量，适合检索、分类和特征提取这种不需要生成文本的场景。CLIP 就是这类模型的招牌，它在 ImageNet 上的零样本分类能力已经能和有监督训练的 ResNet-50 差不多。生成型 VLM 带自回归文本解码器，可以根据视觉输入一个 token 一个 token 地吐出自然语言，适合图像描述、视觉问答和指令跟随这类任务。BLIP-2 的架构亮点是搞了个 Q-Former 当视觉和语言之间的轻量查询桥接器，把视觉特征压缩成少量查询 token 再注入冻结的 LLM，不用重新训整个 LLM，省成本。LLaVA 的路子更直接——一个线性投影就把视觉 token 映射到语言模型的 embedding 空间，然后用视觉指令微调数据让模型学会对话和多轮交互。Qwen-VL 在它们基础上又优化了视觉编码器和语言模型的原生结合，支持任意分辨率的图像输入和更长的视觉 token 序列。

我们这回的设计在语义识别任务上用了对齐型 VLM（CLIP），拿它的零样本分类能力在 40 类的有限空间里做 EEG 辅助分类；在生成式描述任务上用了生成型 VLM（Qwen2-VL-2B-Instruct），靠它的视觉-语言对齐能力和自回归生成能力来生成自由形式的图像描述。

\placeholderfigure{CLIP 与生成式视觉语言模型结构区别示意图}{clip-vlm-contrast-placeholder}{3.2cm}

\subsection{CLIP 图文语义空间}

CLIP 用对称的 InfoNCE 对比损失来训练。假设一个批次里有 $N$ 对图文数据 $\{(x_i, t_i)\}_{i=1}^{N}$，图像编码器输出 $f_{\mathrm{img}}(x_i)$，文本编码器输出 $f_{\mathrm{text}}(t_i)$。训练就是要让正确配对的余弦相似度尽量大，错误配对的相似度尽量小：

\begin{equation}
  \mathcal{L} = -\frac{1}{2N}\sum_{i=1}^{N}\left[
    \log\frac{\exp(\tau^{-1}\cos(f_{\mathrm{img}}(x_i), f_{\mathrm{text}}(t_i)))}
         {\sum_{j=1}^{N}\exp(\tau^{-1}\cos(f_{\mathrm{img}}(x_i), f_{\mathrm{text}}(t_j)))} +
    \log\frac{\exp(\tau^{-1}\cos(f_{\mathrm{img}}(x_i), f_{\mathrm{text}}(t_i)))}
         {\sum_{j=1}^{N}\exp(\tau^{-1}\cos(f_{\mathrm{img}}(x_j), f_{\mathrm{text}}(t_i)))}
  \right]
\end{equation}

这里的 $\tau$ 是一个可训练的温度参数，控制相似度分布的集中程度。

CLIP 有个关键本事是零样本分类。给定 $K$ 个类别，对每个类别 $c$ 造一条文本提示比如 ``a photo of a \{class\_name\}''，用 CLIP 文本编码器算出一个文本原型 $p_c = f_{\mathrm{text}}(t_c)$。对于要分类的图像 $x$，预测结果就是和图像嵌入余弦相似度最高的那个原型对应的类别：

\begin{equation}
  \hat{c} = \arg\max_{c\in\{1,\dots,K\}} \frac{f_{\mathrm{img}}(x)^{\top} p_c}{\|f_{\mathrm{img}}(x)\|\,\|p_c\|}
\end{equation}

我们的设计充分利用了这个特性：用 CLIP text encoder 提前算好 40 个类别的文本原型，训练和评估的时候都冻住不动；模型只需要学会把融合后的多模态嵌入往正确的原型方向推就行，不用在那么点 EEG 数据上从头训一个分类器。

\subsection{EEG 视觉语义解码}

EEG 是通过在头皮上贴电极阵列来记录大脑皮层电活动的一种手段。在视觉认知实验范式里，EEG 能抓住被试看图时产生的诱发响应，像早期的 P100/N170 视觉成分和晚期的 P300 这种语义加工成分。EEG 最大的好处是时间分辨率能到毫秒级，特别适合跟踪快速的神经编码过程。

不过用 EEG 做视觉语义解码有不少固有麻烦：(1) 信噪比低——单次 trial 的信号很容易被肌电、眼动和环境电磁干扰盖住；(2) 个体差异大——不同被试因为电极位置和颅骨厚度不一样，响应模式差异很明显；(3) trial 级数据量小——典型实验里每个被试在每个类别上能用的 trial 就几十条。这些因素决定了直接从少量 EEG 样本学自由文本生成这条路有严重的数据瓶颈。

针对上面这些难题，已经有不少研究者做了各种尝试。Spampinato 他们最早用深度学习方法从 EEG 里解码视觉内容，在物体分类任务上做出了重要结果。Palazzo 他们搞了 EEG-ImageNet 这个大数据集，把 ImageNet 图像和同步采集的 EEG 配对起来。最近几年又有工作进一步探索了把 EEG 嵌入对齐到 CLIP 语义空间，以及用 EEG 做条件生成图像。不过现有研究大多是在完整视觉条件下做的，视觉退化了以后 EEG 到底能起多大语义辅助作用，还没有被系统性地对照验证过。

我们在这个设计里用了这么几条数据利用策略：在 CLIP 预训练好的语义空间里做 EEG-image 对齐，这样就能借助 CLIP 在大规模数据上学到的语义结构，减少对 EEG 数据量的依赖；同时设了 shuffled EEG（打乱脑电）和 random EEG 两种对照组，确保性能提升确实来自真实配对的 EEG 语义信息，而不是别的啥因素。

\subsection{视觉退化建模}

为了系统地评估模型在不同图像质量下的表现，我们定义了六种视觉退化条件：

\begin{itemize}
  \item \textbf{clean}：原始图像，啥处理也不做；
  \item \textbf{lowres16}：图像缩到 $16\times16$ 像素再放大回 $224\times224$，模拟极低分辨率的情况；
  \item \textbf{strong\_blur}：加 $\sigma=5$ 的高斯模糊，模拟运动模糊或者对焦不准；
  \item \textbf{strong\_noise}：加标准差 0.15 的高斯噪声，模拟传感器噪声；
  \item \textbf{occlusion50}：随机遮掉 50\% 的图像区域（Cutout），模拟部分被挡住；
  \item \textbf{mixed}：随机把两种以上退化类型混在一起，模拟复杂的复合退化。
\end{itemize}

这些退化操作在图像喂进 CLIP 编码器之前在线执行，不用每种退化都单独存一份完整图像，存储开销小很多。实验里我们预期真实 EEG 在退化严重的时候增益更大，因为这时候视觉编码器输出不靠谱了，模型更得靠 EEG 辅助信息来补救。

\subsection{多模态融合机制}

我们的设计涉及从粗到细的多层级融合方法：

Embedding 级融合就是把不同模态的向量表示在特征空间里组合起来，包括直接拼接、加权求和和门控残差融合。我们用的门控残差融合公式是 $f_{\mathrm{fused}} = f_{\mathrm{img}} + \sigma(g([f_{\mathrm{img}}; f_{\mathrm{eeg}}])) \odot \Delta(f_{\mathrm{eeg}})$，其中 $g(\cdot)$ 是门控网络，$\sigma$ 是 sigmoid 函数，$\Delta(\cdot)$ 是 EEG 残差变换网络。门控机制让模型能根据联合输入动态调节 EEG 信号的注入强度——图像质量好的时候少加点 EEG，退化严重的时候就多让 EEG 来修正。

Token 级融合是在 CLIP ViT 输出的 visual patch token 序列上插入 EEG token，用 cross-attention 实现：visual token 当 Query，EEG token 当 Key 和 Value，输出增强后的视觉 token。和 embedding 级融合比，token 级融合保留了视觉信息的空间结构，更适合生成式 VLM 的输入格式。

\subsection{LoRA 参数高效微调}

大语言模型和视觉语言模型通常参数多到几十亿，在小规模训练数据上做全参数微调很容易过拟合，还把原来学好的能力搞丢了。LoRA（Low-Rank Adaptation）通过在注意力层的权重矩阵里加一些可训练的低秩分解矩阵来解决这个问题。

对于预训练权重矩阵 $W_0 \in \mathbb{R}^{d\times k}$，LoRA 注入一个低秩增量 $\Delta W = BA$，其中 $B \in \mathbb{R}^{d\times r}$，$A \in \mathbb{R}^{r\times k}$，秩 $r$ 比 $d$ 和 $k$ 小得多。前向计算变成 $h = W_0 x + \frac{\alpha}{r} \cdot BA x$，$\alpha$ 是缩放因子。训练的时候只更新 $A$ 和 $B$，原始权重 $W_0$ 一直冻着。

在我们这个课程设计的生成式描述任务里，LoRA 加在 Qwen2-VL 注意力机制的 Query 和 Value 投影矩阵上，秩 $r=8$。用 LoRA 而不是全参数微调，是我们能在单 GPU 上跑通 Qwen2-VL-2B 生成模型训练的关键工程手段。
